\documentclass[11pt]{article}
\usepackage[T1]{fontenc}
\usepackage{textcomp}
\usepackage[margin=0.75in]{geometry}
\usepackage{amsmath,amssymb,amsthm}
\usepackage{array,booktabs}
\usepackage{hyperref}
\setlength{\parindent}{0pt}
\newtheorem{theorem}{Theorem}
\theoremstyle{definition}
\newtheorem{definition}{Definition}
\title{Flow Proof Artifact: Group-cancel-left}
\author{Generated by \texttt{flow doc proof}}
\date{Flow Proof Book}
\begin{document}
\maketitle
Left cancellation holds in a group.\\[0.5em]
\textbf{Source.} dummit-foote — *Abstract Algebra*, §1.1\\[0.5em]
\bigskip
\noindent{\large \textbf{Derived fact 1.} \textit{a * b = a * c implies b = c} \hfill Group \cdot multiplication \cdot left cancellation holds}\par
\smallskip\noindent\textit{Coordinate: Group \cdot multiplication \cdot left cancellation holds}\par
\smallskip\noindent\textit{Source: dummit-foote}\par
\smallskip\noindent\textit{Built on: left inverse recovers the identity, for inverse on Group, parentheses do not matter, for multiplication on Group, one is the left identity, for identity on Group}\par
\medskip
\noindent\textbf{Goal.} a * b = a * c implies b = c\par
\noindent$\displaystyle \forall a \in Group \forall b \in Group \forall c \in Group\quad b = c$\par
\medskip
\noindent
\begin{tabular}{@{} >{\raggedright\arraybackslash}p{0.06\textwidth} >{\raggedright\arraybackslash}p{0.47\textwidth} >{\raggedleft\arraybackslash}p{0.41\textwidth} @{}}
\textbf{\#} & \textbf{Proof} & \textbf{Mathematics} \\
\midrule
\textbf{1.} & We prove that left cancellation holds for multiplication on Group. & \\[0.45em]
\textbf{2.} & We split into exhaustive cases — the claim must hold in each one. & \\[0.45em]
\textbf{3.} & Case 1 (see step 2): suppose a * b  equals  a * c. & \\[0.45em]
\textbf{4.} & We invoke the definitional clause governing inverse on Group: left inverse recovers the identity, for inverse on Group (instantiated for a). & \\[0.45em]
\textbf{5.} & We invoke the definitional clause governing multiplication on Group: parentheses do not matter, for multiplication on Group (instantiated for inv(a), a, b). & \\[0.45em]
\textbf{6.} & We invoke the definitional clause governing identity on Group: one is the left identity, for identity on Group (instantiated for b). & \\[0.45em]
\textbf{7.} & From step 3, step 4, step 5, and step 6, this implies b equals c. Together with the other cases (step 3), the goal is discharged. Hence proven. & $\displaystyle b = c$ \\[0.45em]
\bottomrule
\end{tabular}
\label{thm:Group-multiplication-left_cancellation_holds}
\par\medskip\hrule
\end{document}
